\documentclass[a4j,dvipdfmx]{jsarticle}
\usepackage{amsmath,amssymb}
\usepackage{siunitx}
\usepackage{ascmac}
\usepackage{amsthm}
\usepackage{bm}

\usepackage[margin=10truemm]{geometry}

\newcommand{\R}{\mathbb{R}}
\newcommand{\rank}{\mathrm{rank}}
\newcommand{\tr}{\mathrm{tr}}
\newcommand{\Ker}{\mathrm{Ker}\hspace{0.5mm}}
\newcommand{\E}{\mathcal{E}}
\newcommand{\F}{\mathcal{F}}
\renewcommand{\labelenumi}{(\arabic{enumi})}
\renewcommand{\Im}{\mathrm{Im}\hspace{0.5mm}}
% \usepackage[dvipdfmx]{hyperref}
\begin{document}
    \part*{第五回 線型代数 問題}
    \subsection*{問題1}
        $n$次正方行列$A$に対し, $|A|$が$A$の固有値の積であり, $\tr A$が固有値の和であることを示せ.
    \subsection*{問題2}
        次の行列をユニタリ行列(もしくは直交行列)で対角化せよ.
        \begin{equation*}
            (1)\begin{pmatrix}
                1 & i\\i & 1
            \end{pmatrix}\quad
            (2)\begin{pmatrix}
                2 & 1 & 1\\
                1 & 2 & 1\\
                1 & 1 & 2
            \end{pmatrix}
        \end{equation*}
    \subsection*{問題3}
        次の行列が対角化可能なら対角化し, そうでなければ理由を述べよ.
        \begin{equation*}
            (1)\begin{pmatrix}
                6 & -3 & -8\\
                -1 & 2 & 1\\
                5 & -3 & -6
            \end{pmatrix}\quad 
            (2)\begin{pmatrix}
                1  & 2  & 1 \\
                -1 & 4  & 1 \\
                2  & -4 & 0 
            \end{pmatrix}\quad
            (3)\begin{pmatrix}
                1  & 3  & -2 \\
                -3 & 13  & -7 \\
                -5  & 19 & -10 
            \end{pmatrix}
        \end{equation*}
    \subsection*{問題4}
        次の問に答えよ. 
        \begin{enumerate}
            \item $|\bm{x}|=1$のとき, 二次形式${}^t \bm{x}A\bm{x}$の最大値は$A$の固有値の最大値であり, 最小値は$A$の固有値の最小値であることを示せ.
            \item $x^2+y^2+z^2=1$のとき, $f(x,y,z)=5x^2+2y^2+5z^2+4xy+4yz-2zx$の最大値と最小値を求めよ.
            \item (2)で求めた最大値, 最小値をとる$(x,y,z)$の例を一つずつ求めよ.
        \end{enumerate}
    \subsection*{問題5}
        任意の対称行列は直交行列で対角化可能であることを示せ.
    \subsection*{問題6}
        3次実行列$A$の固有値が$1,2,3$であるとき, $A^3$を$A^2,A,E$の一次結合でかけ.
    \subsection*{問題7}
        多変数解析学によれば, 二回連続微分可能な二変数関数$f(x,y)$は, 各点に対しある$0<\theta<1$があって, 次のように書ける.
        \begin{equation*}
            f(x+\Delta x,y+\Delta y) = f(x,y)+\left(\Delta x\frac{\partial}{\partial x}+\Delta y\frac{\partial}{\partial y}\right)f(x,y)+\left(\Delta x\frac{\partial}{\partial x}+\Delta y\frac{\partial}{\partial y}\right)^2 f(x+\theta \Delta x,y+\theta \Delta y)\text{\hspace{3mm}(\textbf{Taylor})}
        \end{equation*}
        この式を用いて, 極値を取りうる点でHessianが正なら極値点となり, 負なら鞍点となることを示せ. また, 特に極値を取るときの$f_{xx}$の符号と極大値, 極小値の関係を示せ.
    \subsection*{補充問題}
        $W_1,W_2$をベクトル空間$V$の部分空間とする. 以下の\textbf{次元公式}を次の手順で示せ.
        \begin{equation*}
            \dim(W_1+W_2)=\dim W_1+\dim W_2 - \dim(W_1\cap W_2)
        \end{equation*}
        \begin{enumerate}
            \item $\dim W_1=a, \dim W_2=b ,\dim(W_1\cap W_2)=c$とおく. $W_1\cap W_2$の基底を$W_1,W_2$にそれぞれ拡大して, $W_1,W_2$の基底が作れることを示せ. (一言でもよい)
            \item 得られた$W_1,W_2$の基底のうち, $W_1\cap W_2$に含まれないベクトルの数を$a,b,c$の中から必要なものを使って答えよ.
            \item 得られた$W_1,W_2$の基底のうち, $W_1\cap W_2$に含まれないベクトルと, $W_1\cap W_2$の基底を合わせたものは, $W_1+W_2$の基底であることを示せ.
        \end{enumerate}
\end{document}